\documentclass[11pt,a4paper]{article}
\usepackage{webclases-lesson}
\lessonSetup{T1}{v0.4}{T1 - Movimiento rectilineo}
\begin{document}
\shorthandoff{<>}
\color{Ink}
\coverHeader{T1 --- Movimiento rectil\'ineo}{Una coordenada, dos f\'ormulas y todas las preguntas que se les pueden hacer.}
La cinem\'atica empieza eligiendo un eje, un origen y un reloj. En esta lecci\'on el
movimiento ocurre sobre una recta, as\'i que basta un solo n\'umero para saber d\'onde est\'a
el objeto. Con las integrales de T0 obtendremos dos f\'ormulas de la posici\'on, los
\textbf{santos griales}, y de ellas sacaremos todo lo dem\'as.
\begin{lessonindex}
\indexentry{sec:b1}{1. Una sola coordenada}
\indexentry{sec:b2}{2. Los santos griales: de la aceleraci\'on a la posici\'on}
\indexentry{sec:b3}{3. Qu\'e se le puede preguntar a la f\'ormula}
\indexentry{sec:b4}{4. Lo que cuentan las gr\'aficas}
\end{lessonindex}
\precisionnote

\block{1}{Una sola coordenada}{sec:b1}
Un movimiento es \textbf{rectil\'ineo} o \textbf{unidimensional} cuando el objeto se mueve
siempre sobre la misma recta. Sobre esa recta elegimos un eje: un origen $x=0$, un sentido
positivo y un instante $t=0$. La elecci\'on es libre (el origen puede estar en la salida,
en el suelo o en una farola), pero una vez hecha se mantiene en todo el problema.

\begin{keybox}
\textbf{Por qu\'e aqu\'i basta una f\'ormula.} En general la posici\'on es un vector con
varias componentes. Si el movimiento ocurre sobre una recta y ponemos el eje sobre ella,
todas las componentes salvo una valen siempre cero: hacemos una \textbf{reducci\'on
dimensional}. Solo queda una magnitud, la posici\'on $x$ (o $y$ si el eje es vertical), y su
evoluci\'on se describe con una \'unica funci\'on $x(t)$. En T2 habr\'a una funci\'on por componente.
\end{keybox}

La posici\'on $x$ lleva signo: puede ser negativa, y la distancia al origen es $|x|$.
El desplazamiento $\Delta x=x_f-x_i$ depende solo de los extremos.
La distancia recorrida $d$ suma las longitudes de todos los tramos.

\begin{examplebox}{ir y volver no es quedarse quieto}
\origin{Elaboraci\'on propia.}
\step{Planteamiento} Un alumno sale de $x=2,00\units{m}$, llega a $8,00\units{m}$
y vuelve a $5,00\units{m}$ en un total de $6,00\units{s}$. Compara desplazamiento,
distancia, velocidad media y rapidez media.
\solutionblock
\step{Modelo} Conservamos el mismo eje en ambos tramos.
\begin{center}
\begin{tikzpicture}[x=.8cm,y=1cm,font=\small]
\draw[axis] (0,0)--(9,0) node[right] {$x$ (m)};
\foreach \x in {0,2,5,8} \draw (\x,.06)--(\x,-.06) node[below] {$\x$};
\draw[vector] (2,.4)--(8,.4) node[midway,above] {ida};
\draw[vectoramber] (8,1.1)--(5,1.1) node[midway,above] {vuelta};
\end{tikzpicture}
\end{center}
\step{C\'alculo} $\Delta x=5-2=+3,00\units{m}$;
$d=|8-2|+|5-8|=9,00\units{m}$.
\[
v_{\mathrm{med}}=\frac{\Delta x}{\Delta t}=+0,500\mps,\qquad
\text{rapidez media}=\frac d{\Delta t}=1,50\mps.
\]
\step{Interpretaci\'on} La primera indica el avance neto; la segunda, el ritmo de
recorrido. Si terminase en $x=2\units{m}$, su velocidad media ser\'ia cero,
aunque hubiese caminado.
\end{examplebox}

\heading{Velocidad y aceleraci\'on: dos derivadas}
Como en T0, la velocidad es la derivada de la posici\'on y la aceleraci\'on, la derivada
de la velocidad:
\[
v(t)=\frac{dx}{dt},\qquad a(t)=\frac{dv}{dt}.
\]
$v$ es la pendiente de $x(t)$ en cada instante, no el cociente $x/t$.
La rapidez es $|v|$. Los signos de $v$ y $a$ dicen qu\'e ocurre:
\par\medskip
\begin{tabularx}{\linewidth}{@{}p{2.5cm}Y Y@{}}
\toprule
Signos & Qu\'e hace el m\'ovil & Qu\'e hace su rapidez\\
\midrule
$v>0,\ a>0$ & Avanza hacia $+x$ & Aumenta\\
$v>0,\ a<0$ & Avanza hacia $+x$ & Disminuye\\
$v<0,\ a>0$ & Avanza hacia $-x$ & Disminuye\\
$v<0,\ a<0$ & Avanza hacia $-x$ & Aumenta\\
\bottomrule
\end{tabularx}
\par\medskip
Una pelota lanzada hacia arriba que despu\'es vuelve a caer tambi\'en es un movimiento
unidimensional: se mueve sobre una recta vertical. Con $+y$ hacia arriba, su aceleraci\'on
es negativa todo el rato, $a=-g$. Al subir, $v>0$ y $a<0$: pierde rapidez. Al bajar,
$v<0$ y $a<0$: gana rapidez. La resolveremos entera en el bloque 3.

\block{2}{Los santos griales: de la aceleraci\'on a la posici\'on}{sec:b2}
Si conocemos la aceleraci\'on, \textbf{integrando} obtenemos la velocidad, e
integrando otra vez, la posici\'on. Cada integral deja una constante, y esa constante es
un dato inicial: la posici\'on o la velocidad en $t=0$. Tomamos $t=0$ al empezar a contar.

\heading{Velocidad constante (MRU): una integral}
En el movimiento rectil\'ineo uniforme, $a=0$ y la velocidad $v$ no cambia. Integramos
una vez:
\[
x(t)=x_0+\int_0^t v\,d\tau=x_0+vt.
\]
Para que la f\'ormula quede totalmente definida hacen falta \textbf{dos datos}:
la posici\'on inicial $x_0$ y la velocidad $v$.
La integral es el \'area del rect\'angulo de altura $v$ y base $t$ bajo $v(t)$.

\heading{Aceleraci\'on constante (MRUA): dos integrales}
En el movimiento rectil\'ineo uniformemente acelerado, $a$ es constante. La primera integral
da la velocidad y la segunda, la posici\'on:
\[
v(t)=v_0+\int_0^t a\,d\tau=v_0+at,\qquad
x(t)=x_0+\int_0^t (v_0+a\tau)\,d\tau=x_0+v_0t+\tfrac12at^2.
\]
Ahora hacen falta \textbf{tres datos}: la posici\'on inicial $x_0$, la velocidad inicial $v_0$
y la aceleraci\'on $a$. La segunda integral es el \'area del trapecio bajo $v(t)$:
un rect\'angulo $v_0t$ m\'as un tri\'angulo $\tfrac12at^2$.

\begin{formulabox}
\textbf{Los santos griales: la posici\'on en funci\'on del tiempo.}
\[
\text{MRU:}\quad x(t)=x_0+vt
\qquad\qquad
\text{MRUA:}\quad x(t)=x_0+v_0t+\tfrac12at^2
\]
La deducci\'on con integrales explica de d\'onde salen. En el examen no hace falta
deducirlas: basta con recordarlas y saber sacarles todo lo que contienen.
\end{formulabox}

\heading{La velocidad no se memoriza: se deriva}
Cada vez que necesitemos la velocidad en un MRUA, derivamos el santo grial:
\[
v(t)=\frac{dx}{dt}=\frac{d}{dt}\Big(x_0+v_0t+\tfrac12at^2\Big)=v_0+at,
\qquad a(t)=\frac{dv}{dt}=a.
\]
As\'i no hace falta recordar una segunda f\'ormula: sale en una l\'inea con la regla de
la potencia de T0. En un MRU, derivar $x_0+vt$ da la velocidad constante $v$.
Si el reloj empieza en $t_0\ne0$, se escribe $t-t_0$ en lugar de $t$.

\heading{Qu\'e significa cada coeficiente}
El santo grial del MRUA es un polinomio de segundo grado en $t$. Cada coeficiente tiene
un significado f\'isico y uno gr\'afico:
\par\medskip
\begin{tabularx}{\linewidth}{@{}p{1.9cm}p{3.9cm}Y@{}}
\toprule
T\'ermino & Significa & En la gr\'afica $x(t)$\\
\midrule
$x_0$ (sin $t$) & D\'onde est\'a en $t=0$ & Donde la curva corta el eje vertical\\
$v_0\,t$ & Velocidad inicial & Pendiente de la curva en $t=0$; su signo da el sentido inicial\\
$\tfrac12a\,t^2$ & Mitad de la aceleraci\'on & Curvatura: $a>0$ abre hacia arriba ($\cup$), $a<0$ hacia abajo ($\cap$)\\
\bottomrule
\end{tabularx}
\par\medskip
\textbf{Trampa t\'ipica:} el n\'umero que acompa\~na a $t^2$ es $a/2$, no $a$.
En una recta (MRU) no hay t\'ermino en $t^2$: la gr\'afica no se curva.

\begin{exercisebox}{leer un santo grial}
Un m\'ovil cumple $x(t)=4+8t-2t^2$ en SI. Identifica $x_0$, $v_0$ y $a$, escribe
$v(t)$ derivando y halla d\'onde da la vuelta.
\solutionblock
$x_0=4\units{m}$, $v_0=8\mps$ y $\tfrac12a=-2$, luego $a=-4\mpss$.
$v=dx/dt=8-4t$, que se anula en $t=2,00\units{s}$. Ah\'i $x(2)=4+16-8=12,0\units{m}$:
es el punto m\'as alejado, porque $a<0$ y la par\'abola abre hacia abajo.
\end{exercisebox}

\block{3}{Qu\'e se le puede preguntar a la f\'ormula}{sec:b3}
Con el santo grial escrito, casi todas las preguntas de un problema son una de estas.
Solo las dos marcadas con $\bigstar$ tienen algo de dificultad.
\par\medskip
\begin{tabularx}{\linewidth}{@{}Y Y@{}}
\toprule
Pregunta & C\'omo se responde\\
\midrule
\textquestiondown D\'onde est\'a en el instante $t_1$? & Sustituir: $x(t_1)$\\
\textquestiondown Qu\'e velocidad lleva en $t_1$? & Derivar y sustituir: $v(t_1)$\\
\textquestiondown Cu\'ando pasa por la posici\'on $x_1$? & Resolver $x(t)=x_1$: primer grado en MRU, segundo en MRUA\\
$\bigstar$ \textquestiondown Cu\'ando se para? \textquestiondown Altura m\'axima o posici\'on m\'inima? & Imponer $v(t)=0$; despu\'es sustituir ese $t$ en $x(t)$\\
$\bigstar$ \textquestiondown Cu\'ando llega al suelo y con qu\'e velocidad? & Resolver $x(t)=0$ (o $y=0$), un polinomio de grado 2; elegir la ra\'iz y sustituirla en $v(t)$\\
\textquestiondown Cu\'ando y d\'onde se encuentran dos m\'oviles? & Igualar sus santos griales: $x_A(t)=x_B(t)$\\
\bottomrule
\end{tabularx}
\par\medskip
\textbf{Parar no es llegar al origen.} El punto de giro es donde $v=0$, que es justo
donde $x(t)$ tiene un m\'aximo o un m\'inimo, como en T0. Llegar al suelo o al origen es
$x=0$, otra condici\'on distinta.

\begin{toolbox}[breakable]{resolver el polinomio de segundo grado}
\toolpart{Qu\'e hace} Encuentra los instantes en que un MRUA pasa por una posici\'on dada.
\toolpart{C\'omo} Pasa todo a un lado y ordena: $At^2+Bt+C=0$ (en may\'usculas, para no
confundir $A$ con la aceleraci\'on). Entonces
\[
t=\frac{-B\pm\sqrt{B^2-4AC}}{2A}.
\]
\toolpart{Qu\'e dice el discriminante} Si $B^2-4AC<0$, nunca llega a esa posici\'on
(por ejemplo, la altura pedida est\'a por encima del m\'aximo). Si vale $0$, la toca justo
una vez: es el punto de giro. Si es positivo, hay dos instantes.
\toolpart{Elegir la ra\'iz} Descarta los $t<0$ (son antes de empezar) y los que caen fuera
del tramo descrito. Si quedan dos positivas, las dos pueden ser f\'isicas: una al subir y
otra al bajar. Comprueba sustituyendo en $x(t)$.
\toolpart{Atajo} Si $C=0$, saca factor com\'un: $t(At+B)=0$ da $t=0$ y $t=-B/A$.
\end{toolbox}

\begin{examplebox}{la pelota que sube y vuelve a caer}
\origin{Adaptaci\'on del tiro vertical del original T1 con datos propios.}
\step{Planteamiento} Una pelota sale hacia arriba a $18,0\mps$ desde la mano, a
$1,50\units{m}$ del suelo. Sin aire y con $g=10\mpss$, calcula la altura m\'axima, cu\'ando
llega al suelo y con qu\'e velocidad.
\solutionblock
\step{F\'isicamente} La pelota sube cada vez m\'as despacio, se para un instante arriba y
baja cada vez m\'as deprisa. Todo ocurre sobre la misma vertical: es unidimensional.
La gravedad tira hacia abajo durante todo el vuelo, tambi\'en arriba: la aceleraci\'on no
cambia de sentido.
\step{Matem\'aticamente} Eje $+y$ hacia arriba con origen en el suelo: $y_0=1,5\units{m}$,
$v_0=+18\mps$ y $a=-10\mpss$. El santo grial y su derivada son
\[
y(t)=1,5+18t-5t^2,\qquad v(t)=\frac{dy}{dt}=18-10t.
\]
El coeficiente $-5$ es $a/2$.
\step{Altura m\'axima ($\bigstar$)} $v=0$ da $18-10t=0$, $t=1,80\units{s}$, y
$y(1,8)=1,5+32,4-16,2=17,7\units{m}$. Sobre la mano ha subido $16,2\units{m}$.
\step{Llegada al suelo ($\bigstar$)} $y=0$ da $-5t^2+18t+1,5=0$, con
$A=-5$, $B=18$ y $C=1,5$:
\[
t=\frac{-18\pm\sqrt{18^2-4(-5)(1,5)}}{2(-5)}=\frac{-18\pm\sqrt{354}}{-10}
\quad\Rightarrow\quad t=3,68\units{s}\ \text{o}\ t=-0,0815\units{s}.
\]
La ra\'iz negativa es anterior al lanzamiento y se descarta. Entonces
$v(3,68)=18-36,8=-18,8\mps$: el signo menos indica que baja.
\step{Gr\'aficamente} Las tres gr\'aficas comparten reloj. La l\'inea discontinua marca
$t=1,80\units{s}$: m\'aximo de $y$, cero de $v$, y la aceleraci\'on sigue valiendo $-10\mpss$.
\begin{center}
\begin{tikzpicture}[font=\small]
\begin{scope}[x=.9cm,y=.1cm]
\draw[axis] (0,0)--(4.2,0) node[right] {$t$ (s)};
\draw[axis] (0,0)--(0,21) node[above] {$y$ (m)};
\draw[BrandBlue,thick,domain=0:3.6815,samples=60] plot(\x,{1.5+18*\x-5*\x*\x});
\draw[guide] (1.8,0)--(1.8,17.7)--(0,17.7);
\node[left] at (0,17.7) {$17,7$}; \node[left] at (0,1.5) {$1,5$};
\node[below] at (1.8,0) {$1,8$}; \node[below] at (3.68,0) {$3,68$};
\end{scope}
\begin{scope}[xshift=5.6cm,yshift=1.05cm,x=.9cm,y=.05cm]
\draw[axis] (0,0)--(4.2,0) node[right] {$t$ (s)};
\draw[axis] (0,-21)--(0,21) node[above] {$v$ (m/s)};
\draw[BrandBlue,thick] (0,18)--(3.6815,-18.815);
\draw[guide] (1.8,-21)--(1.8,21);
\node[left] at (0,18) {$18$}; \node[right] at (3.68,-18.8) {$-18,8$};
\node[above right] at (1.8,0) {$1,8$};
\end{scope}
\begin{scope}[xshift=11.1cm,yshift=1.95cm,x=.9cm,y=.13cm]
\draw[axis] (0,0)--(4.2,0) node[right] {$t$ (s)};
\draw[axis] (0,-12)--(0,4) node[above] {$a$ (m/s$^2$)};
\draw[BrandBlue,thick] (0,-10)--(3.6815,-10);
\draw[guide] (1.8,-12)--(1.8,4);
\node[left] at (0,-10) {$-10$};
\end{scope}
\end{tikzpicture}
\end{center}
$y(t)$ es una par\'abola que abre hacia abajo porque $a<0$; $v(t)$ es una recta de
pendiente $-10$; $a(t)$ es una recta horizontal.
\step{Interpretaci\'on} En lo m\'as alto $v=0$, pero $a\ne0$: si la aceleraci\'on fuese cero
all\'i, la pelota se quedar\'ia flotando. A los $3,60\units{s}$ vuelve a la altura de la mano
con $-18,0\mps$, la misma rapidez con la que sali\'o. Llega al suelo algo despu\'es y m\'as
r\'apida porque cae $1,50\units{m}$ m\'as.
\end{examplebox}

\begin{examplebox}{encuentro en el mismo eje}
\origin{Elaboraci\'on propia, con el m\'etodo del encuentro del original T1.}
\step{Planteamiento} A sale de $-20,0\units{m}$ a $+2,00\mps$ y B de
$40,0\units{m}$ a $-1,00\mps$, simult\'aneamente. Halla cu\'ando y d\'onde coinciden.
\solutionblock
\step{Modelo} Ambos llevan MRU, as\'i que cada uno tiene su santo grial:
\[
x_A=-20+2t,\qquad x_B=40-t,\qquad x_A(t)=x_B(t).
\]
\step{C\'alculo} $3t=60$ da $t=20,0\units{s}$.
Sustituir en cualquiera de las dos leyes da $x=20,0\units{m}$.
\begin{center}
\begin{tikzpicture}[x=.18cm,y=.047cm,font=\small]
\draw[axis] (0,0)--(34,0) node[right] {$t$ (s)};
\draw[axis] (0,-25)--(0,50) node[above] {$x$ (m)};
\foreach \t in {10,20,30} \draw (\t,0)--(\t,-1.4) node[below] {$\t$};
\foreach \p in {-20,20,40} \draw (0,\p)--(-.5,\p) node[left] {$\p$};
\draw[BrandBlue,thick,domain=0:30] plot(\x,{-20+2*\x});
\draw[GreenLine!70!black,dashed,thick,domain=0:30] plot(\x,{40-\x});
\fill (20,20) circle[x radius=.5,y radius=1.8];
\draw[guide] (20,0)--(20,20);
\node[right] at (30,40) {A};
\node[right] at (30,10) {B};
\end{tikzpicture}
\end{center}
\step{Interpretaci\'on} A recorre $40,0\units{m}$ y B $20,0\units{m}$:
coincidir no exige igual distancia recorrida. La intersecci\'on de las rectas
representa una misma posici\'on al mismo tiempo.
\variation{Si B va a $+3,00\mps$, la soluci\'on es $t=-60,0\units{s}$: es anterior a la
salida, as\'i que no hay encuentro.}
\end{examplebox}

\begin{examplebox}{frenar y cambiar de sentido}
\origin{Adaptaci\'on del MRUA del original T1 con los datos de la versi\'on actual.}
\step{Planteamiento} Un carrito parte de $x_0=100\units{m}$ con
$v_0=12,0\mps$ y $a=-3,00\mpss$. Halla instante y posici\'on de parada,
y su estado a los $6,00\units{s}$ si se mantiene la misma aceleraci\'on.
\solutionblock
\step{Modelo} Santo grial y derivada:
\[
x(t)=100+12t-1,5t^2,\qquad v(t)=\frac{dx}{dt}=12-3t.
\]
Inicialmente $v$ y $a$ son opuestos: pierde rapidez. La parada es $v=0$, no $x=0$.
\step{C\'alculo} $12-3t=0$ da $t=4,00\units{s}$ y $x(4)=100+48-24=124\units{m}$.
A los $6,00\units{s}$: $v(6)=-6,00\mps$ y $x(6)=100+72-54=118\units{m}$.
\step{Interpretaci\'on} Ha avanzado $24,0\units{m}$ y vuelto $6,00\units{m}$:
desplazamiento $18,0\units{m}$, recorrido $30,0\units{m}$.
\variation{Un coche que frena y se queda parado cambia de modelo al detenerse:
despu\'es, $v=0$ y $x=124\units{m}$. Prolongar $a=-3$ inventar\'ia una marcha atr\'as.}
\end{examplebox}

\begin{examplebox}{reacci\'on y frenado}
\origin{Elaboraci\'on propia; modelo ideal por tramos.}
\step{Planteamiento} Un coche va a $72,0\units{km/h}=20,0\mps$.
Mantiene esa velocidad durante $0,800\units{s}$ de reacci\'on y luego
frena con $a=-5,00\mpss$. Halla tiempo y distancia totales de parada.
\solutionblock
\step{Modelo} Dos tramos, dos santos griales. En la reacci\'on hay MRU: $x=20t$, que acaba
en $x=16,0\units{m}$. En la frenada hay MRUA con reloj propio $\tau=t-0,8$, empezando
donde acab\'o el primero:
\[
x=16+20\tau-2,5\tau^2,\qquad v=\frac{dx}{d\tau}=20-5\tau.
\]
\step{C\'alculo} $v=0$ da $\tau=4,00\units{s}$ y $x=16+80-40=56,0\units{m}$.
Total: $t=4,80\units{s}$ y $d=56,0\units{m}$, de los que $40,0\units{m}$ son de frenada.
\variation{A $40,0\mps$, la reacci\'on ocupa $32,0\units{m}$ y la frenada
$160\units{m}$: duplicar la velocidad cuadruplica la frenada.}
\end{examplebox}

\block{4}{Lo que cuentan las gr\'aficas}{sec:b4}
\begin{tabularx}{\linewidth}{@{}p{2.7cm}Y@{}}
\toprule
Representaci\'on & Lectura f\'isica\\
\midrule
$x(t)$ & Pendiente $v$; un m\'aximo o un m\'inimo es un cambio de sentido.\\
$v(t)$ & Pendiente $a$; \'area con signo $\Delta x$.\\
$a(t)$ & \'Area con signo $\Delta v$; hace falta $v_0$ para hallar $v$.\\
\bottomrule
\end{tabularx}
Las \'areas son las mismas integrales del bloque 2:
$x(t)=x_0+\int_0^t v(\tau)\,d\tau$ y $v(t)=v_0+\int_0^t a(\tau)\,d\tau$.
Una pendiente negativa en $x(t)$ significa $v<0$, no necesariamente $a<0$.

\begin{examplebox}{desplazamiento cero, recorrido positivo}
\origin{Elaboraci\'on propia.}
\step{Planteamiento} $v(t)=4-2t$ en SI, $0\le t\le4\units{s}$, $x_0=0$.
Halla cambio de sentido, posici\'on final y recorrido.
\solutionblock
\step{Modelo} El signo de $v$ cambia en $t=2\units{s}$. Separamos ah\'i las \'areas.
\begin{center}
\begin{tikzpicture}[x=1.2cm,y=.35cm,font=\small]
\begin{scope}
\draw[axis] (0,0)--(4.7,0) node[right] {$t$ (s)};
\draw[axis] (0,0)--(0,5.2) node[above] {$x$ (m)};
\draw[BrandBlue,thick,domain=0:4,samples=60] plot(\x,{4*\x-\x*\x});
\draw[guide] (2,0)--(2,4);
\node[left] at (0,4) {$4$}; \node[below] at (2,0) {$2$};
\node[below] at (4,0) {$4$};
\end{scope}
\begin{scope}[xshift=7.2cm]
\fill[GreenLine!20] (0,0)--(0,4)--(2,0)--cycle;
\fill[AmberLine!20] (2,0)--(4,-4)--(4,0)--cycle;
\draw[axis] (0,0)--(4.7,0) node[right] {$t$ (s)};
\draw[axis] (0,-4.5)--(0,5) node[above] {$v$ (m/s)};
\draw[BrandBlue,thick] (0,4)--(4,-4);
\node[left] at (0,4) {$4$}; \node[left] at (0,-4) {$-4$};
\node[above right] at (2,0) {$2$}; \node[above] at (4,0) {$4$};
\node at (.75,1.2) {$+4\units{m}$};
\node at (3.25,-1.2) {$-4\units{m}$};
\end{scope}
\end{tikzpicture}
\end{center}
\step{C\'alculo}
$\Delta x=4-4=0\units{m}$, $d=4+4=8,00\units{m}$.
Integrando, $x(t)=4t-t^2$: es el santo grial con $x_0=0$, $v_0=4\mps$ y $a=-2\mpss$.
El m\'aximo es $x(2)=4,00\units{m}$ y $x(4)=0$.
\step{Interpretaci\'on} Las gr\'aficas comparten reloj: $v=0$ coincide con el
m\'aximo de $x$, no con su paso por el origen. La velocidad media es cero y la rapidez
media $2,00\mps$.
\end{examplebox}

\begin{exercisebox}{qu\'e permite afirmar una tabla}
$t$ en s: $0;\ 1;\ 2;\ 3$.
$x$ en m: $10;\ 12;\ 18;\ 28$.
Bajo la hip\'otesis de MRUA desde reposo, obtiene $a$ y predice $x(4)$.
Justifica si los datos garantizan por s\'i solos ese futuro.
\solutionblock
Con $x=10+\tfrac12at^2$, el primer intervalo da $a=4,00\mpss$.
La misma ley reproduce $x(2)=18$ y $x(3)=28$, y predice
$x(4)=42,0\units{m}$. Una tabla finita es compatible con ese modelo,
pero no demuestra que la aceleraci\'on siga constante fuera de lo observado.
\end{exercisebox}

\heading{Fuentes y procedencia}
Redacci\'on y gr\'aficas propias; MRUA y tiro vertical siguen los m\'etodos del original T1.
\source{https://fisquiweb.es/Apuntes/apun1B.htm}{FisQuiWeb / IES La Magdalena: inspiraci\'on y material original de 1.\textordmasculine.}
\source{https://openstax.org/books/university-physics-volume-1/pages/3-4-motion-with-constant-acceleration}{OpenStax, University Physics 1, 3.4: Motion with Constant Acceleration.}
\end{document}
